\documentclass[10pt,a4paper]{amsart}
\usepackage[margin=19mm]{geometry}
\usepackage{amsmath,amssymb,mathtools}
\usepackage[scr=rsfs,cal=euler]{mathalfa}
\usepackage{xfrac}
\usepackage[unicode,hidelinks,bookmarks=false]{hyperref}
\newcommand{\ie}{i.e.\ }
\newcommand{\cf}{cf.\ }
\newcommand{\resp}{respectively\ }
\newcommand{\Subsection}{\subsection}
\providecommand{\nobreakdash}{\nobreak\hbox{-}}
\setlength{\emergencystretch}{2em}
\multlinegap0pt
\usepackage{mathtools}
\usepackage{algorithm}\usepackage{algpseudocode}
\usepackage[shortlabels]{enumitem}
\usepackage{amssymb}
\usepackage{float}
\usepackage{csquotes}
\newtheorem{definition}{Definition}
\newtheorem{lemma}{Lemma}
\title{Sobolev Norm Learning Rates for Conditional Mean Embeddings}
\author{Prem Talwai, Ali Shameli, David Simchi-Levi}
\date{Source: arXiv:2105.07446v3 (CC BY 4.0)}
\begin{document}
\maketitle

\noindent\emph{Source and licence.} Sobolev Norm Learning Rates for Conditional Mean Embeddings. Version arXiv:2105.07446v3 (CC BY 4.0), distributed by arXiv under the Creative Commons Attribution 4.0 International licence, http://creativecommons.org/licenses/by/4.0/, as recorded in the arXiv licence metadata for this exact version and shown on the abstract page https://arxiv.org/abs/2105.07446v3. Full licence notice: \texttt{licenses/arxiv-2105.07446-CC-BY-4.0.txt}.\par\smallskip

\noindent\textit{Editorial scope.} Counted unit: the article's lemma Spaces (source0.tex lines 224--232). The article's complete original proof of it is delivered with it (source0.tex lines 551--575, one \texttt{proof} environment of 25 lines, which the article places away from the statement). No mathematics is written, repaired or re-derived: the statement and the proof are the article's own lines, in the article's own notation, and no retained line is substituted or re-flowed.  Retained uncounted as the source-local context the proof needs: lines 197--205.  Nothing was omitted from any retained span, so the package carries no omission record.  No retained line contains a citation, so the wrapper carries no bibliography and no bibliography entry.  No retained line contains a figure environment, an image-inclusion command or drawing code, so no image is delivered and the package declares no asset dependency.  Editorial formatting only, in addition: an AMS \texttt{amsart} preamble that restates the article's own declarations and macros used by the retained lines, namely the environments \texttt{definition}, \texttt{lemma} on separate counters, together with the standard packages those lines need.  The retained environments print in this document as Definition 1, Lemma 1 in order of appearance, whereas the article prints the same items with the numbering of its own full document.  The complete native source of the article as distributed in the arXiv e-print for this exact version is bundled unchanged as \texttt{sources/111516/source0.tex}.
\begin{definition}
\label{Sobolev Norm}
Let $T: \mathcal{H}^{\beta}_K \to \mathcal{H}_L$ be a (possibly unbounded) operator for some $\beta > 0$ and let $\gamma \in (0, \beta)$. Let $I_{\beta, \gamma, \nu}: \mathcal{H}^{\beta} \to \mathcal{H}^{\gamma}$ be the canonical embedding. We define the operator norm $||\cdot||_{\gamma}$:
\begin{equation*}
\label{Sobolev Norm Formula}
    ||T||_{\beta, \gamma} = ||T \circ I^{*}_{\beta, \gamma, \nu}||_{\mathcal{H}^{\gamma}_{K} \to \mathcal{H}_L}
\end{equation*}
where both norms may possibly be infinite. When $\beta = 1$, we simply write $||\cdot||_{\gamma}$
\end{definition}

\begin{lemma}
\label{Operator Norms in Interpolation Spaces}
Let $T: \mathcal{H}^{\beta}_K \to \mathcal{H}_L$ be an operator. Then, for any $\gamma \in (0, \beta)$, we have that:
\begin{equation*}
    ||T||_{\beta, \gamma} = ||T \circ C_{\beta, \gamma, \nu}^{\frac{1}{2}}||_{\mathcal{H}^{\beta}_K \to \mathcal{H}_L}
\end{equation*}
where $C_{\beta, \gamma, \nu} = I^{*}_{\beta, \gamma, \nu}I_{\beta, \gamma, \nu}$ 
%(with $I^{*}_{\beta, \gamma, \nu}$ as in Definition \ref{Sobolev Norm}). 
\end{lemma}

\begin{proof}[Proof of Lemma \ref{Operator Norms in Interpolation Spaces}]
Since $\{\mu^{\frac{\beta}{2}}_i e_i\}_{i = 1}^{\infty}$ is an orthonormal basis for $\mathcal{H}^{\beta}_K$, we may express any $f \in \mathcal{H}^{\beta}_{K}$ as $f = \sum_{i = 1}^{\infty} \langle f, \mu^{\frac{\beta}{2}}_i e_i \rangle_{\mathcal{H}^{\beta}_K}\mu^{\frac{\beta}{2}}_i e_i$. Hence, we have:
\begin{align*}
    \langle f, C_{\beta, \gamma, \nu} (\mu^{\frac{\beta}{2}}_i e_i) \rangle_{\mathcal{H}^{\beta}_K} & = \langle f, I^{*}_{\beta, \gamma, \nu}I_{\beta, \gamma, \nu} (\mu^{\frac{\beta}{2}}_i e_i) \rangle_{\mathcal{H}^{\beta}_K} \\
    & = \langle I_{\beta, \gamma, \nu}f, I_{\beta, \gamma, \nu} (\mu^{\frac{\beta}{2}}_i e_i) \rangle_{\mathcal{H}^{\gamma}_K} \\
    & = \langle f, \mu^{\frac{\beta}{2}}_i e_i \rangle_{\mathcal{H}^{\gamma}_K} \\
    & = \Big \langle \sum_{i = 1}^{\infty} \langle f, \mu^{\frac{\beta}{2}}_i e_i \rangle_{\mathcal{H}^{\beta}_K}\mu^{\frac{\beta}{2}}_i e_i, \mu^{\frac{\beta}{2}}_i e_i \Big\rangle_{\mathcal{H}^{\gamma}_K} \\
    & = \mu^{\beta - \gamma}_i \langle f, \mu^{\frac{\beta}{2}}_i e_i \rangle_{\mathcal{H}^{\beta}_K}
\end{align*}
where the final step follows from the fact that $\{\mu_i^{\frac{\gamma}{2}}e_i\}_{i = 1}^{\infty}$ is an orthonormal basis in $\mathcal{H}^{\gamma}_K$. Hence $C_{\beta, \gamma, \nu}$ is a positive self-adjoint operator on $\mathcal{H}^{\beta}_K$ with eigenvalues $\{\mu^{\beta - \gamma}_i\}_{i = 1}^{\infty}$ and an orthonormal basis of eigenfunctions $\{\mu^{\frac{\beta}{2}}_i e_i\}_{i = 1}^{\infty}$. Moreover, since $C_{\beta, \gamma, \nu} = I^{*}_{\beta, \gamma, \nu}I_{\beta, \gamma, \nu}$ by definition and $I_{\beta, \gamma, \nu}$ is the canonical embedding of $H^{\beta}_{K}$ into $H^{\gamma}_{K}$, it follows that the action of $I^{*}_{\beta, \gamma, \nu}: H^{\gamma}_{K} \to  \mathcal{H}^{\beta}_{K}$ can be characterized as:
\begin{equation}
\label{Action of Adjoint}
    I^{*}_{\beta, \gamma, \nu} e_i = \mu^{\beta - \gamma}_{i} e_i \hspace{4mm} \nu-\text{almost surely}
\end{equation}
for all $i \in \mathbb{N}$. Now, let $B_{\ell^2}$ denote the unit ball in $\ell^2$. Then, we have that for any linear operator $T: \mathcal{H}^{\beta}_K \to \mathcal{H}_L$:
\begin{align*}
    ||T||_{\beta, \gamma} & = ||T \circ I^{*}_{\beta, \gamma, \nu}|| \\
    & = \sup_{a \in B_{\ell^2}} \Big\|\sum_{i} a_i \mu^{\frac{\gamma}{2}}_i (T \circ I^{*}_{\beta, \gamma, \nu})  e_i \Big\| \\
    & = \sup_{a \in B_{\ell^2}} \Big\|\sum_{i} \mu^{\beta - \frac{\gamma}{2}}_i a_i (Te_i)\Big\| \hspace{10mm} \text{by \eqref{Action of Adjoint}} \\
    & = \sup_{a \in B_{\ell^2}} \Big\|\sum_{i} a_i\mu^{\frac{\beta - \gamma}{2}}_i  T(\mu^{\frac{\beta}{2}}_ie_i)\Big\| \\
    & = \sup_{a \in B_{\ell^2}} \Big\|\sum_{i} a_i (T \circ C_{\beta, \gamma, \nu}^{\frac{1}{2}})\mu^{\frac{\beta}{2}}_i e_i\Big\| \\
    & = ||T \circ C_{\beta, \gamma, \nu}^{\frac{1}{2}}||
\end{align*}
where again the last and second equalities follow from the fact that $\{\mu^{\frac{\beta}{2}}_i e_i\}_{i \in \mathbb{N}}$ and $\{\mu^{\frac{\gamma}{2}}_i e_i\}_{i \in \mathbb{N}}$ are orthonormal bases in $\mathcal{H}^{\beta}_K$ and $\mathcal{H}^{\gamma}_K$, respectively. 
\end{proof}

\end{document}
